% Related Work — see issue #5.
\section{Related Work}
\label{sec:related_work}

Several previously published studies have informed the development of the proposed
methodology. Collectively, these works represent the current state of the art in
classification techniques for Alzheimer's disease (AD). A comparative overview of
these approaches is provided in \cref{tab:research_matrix}.

% research_matrix.tex — comparison of prior AD-classification studies (issue #5).
% Included from sections/related_work.tex with % research_matrix.tex — comparison of prior AD-classification studies (issue #5).
% Included from sections/related_work.tex with % research_matrix.tex — comparison of prior AD-classification studies (issue #5).
% Included from sections/related_work.tex with \input{../tables/research_matrix}
% (requires the `array` and `booktabs` packages, both loaded in preamble.tex).
%
% Rows were compiled from the performance tables of Nibila & Sivagami (2026),
% so metrics are AS REPORTED BY THAT REVIEW and must be checked against each
% primary paper before the numbers are quoted. `%` comments above each row
% hold the DOI and verification notes; they are not typeset.
\begin{table*}[t]
  \centering
  \footnotesize
  \setlength{\tabcolsep}{3pt}
  \renewcommand{\arraystretch}{1.15}
  \newcolumntype{L}[1]{>{\raggedright\arraybackslash}p{#1}}
  \caption{Summary of previous studies on Alzheimer's disease classification
    from neuroimaging. CN: cognitively normal; MCI: mild cognitive impairment;
    AD: Alzheimer's disease. Metrics are those reported by each study.}
  \label{tab:research_matrix}
  \begin{tabular}{L{2.2cm} L{2.4cm} L{1.5cm} L{2.0cm} L{3.9cm} L{3.7cm}}
    \toprule
    \textbf{Author (year)} & \textbf{Dataset} & \textbf{Modality} & \textbf{Task} &
    \textbf{Architecture} & \textbf{Validation metrics} \\
    \midrule
    % doi: 10.1016/j.neucom.2018.11.111 | note: Unsupervised; dataset not named in review
    Bi et al. (2020) & 1,075 subjects & MRI & AD vs CN & PCANet + K-means (unsupervised) & Acc 92.6\%; Prec 0.84; Rec 0.86 \\
    % doi: 10.1016/j.asoc.2019.105857 | note: Dataset not named in review
    Chitradevi \& Prabha (2020) & 200 subjects & MRI & AD vs CN & Grey Wolf Optimization (hippocampus sub-regions) + AlexNet & Acc 95\%; Prec 0.94; Rec 0.95 \\
    % doi: 10.1016/j.neucom.2020.01.053 | note: Precision/recall exceed accuracy: inconsistent in review
    Li et al. (2020) & ADNI (389 subjects) & 4D fMRI & AD vs CN & 3D CNN & Acc 88.12\%; Prec 0.92; Rec 0.93 \\
    % doi: 10.1016/j.neuroimage.2019.116459 | note: Segmentation + classification pipeline
    Liu et al. (2020) & ADNI (439 subjects) & MRI & CN vs MCI vs AD & Multi-model deep CNN + 3D DenseNet (hippocampus segmentation then classification) & Acc 76.2\%; Prec 0.78; Rec 0.77 \\
    % doi: 10.1016/j.patcog.2021.107944
    Chen \& Xia (2021) & ADNI (1,248 subjects) & MRI & CN vs MCI vs AD & Iterative sparse learning + deep learning & Acc 95.32\%; Prec 0.96; Rec 0.97 \\
    % doi: 10.1016/j.compbiomed.2021.104537
    Ebrahimi et al. (2021) & ADNI (450 subjects) & MRI & CN vs MCI vs AD & Temporal convolutional network (TCN) & Acc 91.78\%; Prec 0.87; Rec 0.90 \\
    % doi: 10.1016/j.bspc.2020.102397
    Hedayati et al. (2021) & ADNI (800 subjects) & MRI & CN vs MCI vs AD & Ensemble of convolutional autoencoders (feature extraction) & Acc 92.5\%; Prec 0.85; Rec 0.88 \\
    % doi: 10.1016/j.irbm.2020.06.006 | note: Modality as reported by review; verify
    Janghel \& Rathore (2021) & 2,729 images & fMRI, PET & AD vs CN & VGG16 + SVM & Acc 76\%; Prec 0.79; Rec 0.80 \\
    % doi: 10.1016/j.compbiomed.2021.104828 | note: Dataset not named in review
    Kaplan et al. (2021) & 263 subjects & MRI & CN vs MCI vs AD & LPQNet & Acc 99\%; Prec 0.93; Rec 0.94 \\
    % doi: 10.1016/j.cmpb.2021.106032 | note: Review and paper title disagree on architecture; review's DOI was wrong
    Liu et al. (2021) & OASIS (416 subjects) & sMRI & AD vs CN & Depthwise separable CNN & Acc 93\%; Prec 0.86; Rec 0.88 \\
    % doi: 10.1016/j.bspc.2021.103455
    Deepa \& Chokkalingam (2022) & ADNI (819 subjects) & MRI & CN vs MCI vs AD & VGG16 + Arithmetic Optimization Algorithm & Acc 97\%; Prec 0.99; Rec 0.98 \\
    % doi: 10.1007/s12559-021-09946-2
    Helaly et al. (2022) & ADNI (300 subjects) & MRI & CN vs MCI vs AD & CNN & Acc 98.47\%; Prec 0.95; Rec 0.96 \\
    % doi: 10.1016/j.bspc.2022.103565 | note: Precision/recall exceed accuracy: inconsistent in review
    Kong et al. (2022) & ADNI (740 subjects) & MRI + PET & AD vs CN & 3D CNN (multimodal) & Acc 86.52\%; Prec 0.94; Rec 0.95 \\
    % doi: 10.1016/j.matpr.2021.04.415 | note: Image-level count; check for subject-level split
    Kumar et al. (2022) & OASIS (5,000 images) & MRI & AD vs CN & AlexNet & Acc 97\%; Prec 0.93; Rec 0.92 \\
    % doi: 10.1016/j.asoc.2021.108099
    Sharma et al. (2022) & ADNI (400 subjects) & sMRI & CN vs MCI vs AD & FDN-ADNet: ResNet + fuzzy LS-TWSVM & Acc 97.29\%; Prec 0.98; Rec 0.98 \\
    % doi: 10.3389/frai.2025.1558725 | note: Public Kaggle slices; task differs between review and title
    Elmotelb et al. (2025) & Kaggle MRI (6,400 images) & MRI & AD vs CN; multi-class & ResNet152V2 + hyperparameter optimization & Acc 95\%; Prec 0.98; Rec 0.98 \\
    \bottomrule
  \end{tabular}
\end{table*}

% (requires the `array` and `booktabs` packages, both loaded in preamble.tex).
%
% Rows were compiled from the performance tables of Nibila & Sivagami (2026),
% so metrics are AS REPORTED BY THAT REVIEW and must be checked against each
% primary paper before the numbers are quoted. `%` comments above each row
% hold the DOI and verification notes; they are not typeset.
\begin{table*}[t]
  \centering
  \footnotesize
  \setlength{\tabcolsep}{3pt}
  \renewcommand{\arraystretch}{1.15}
  \newcolumntype{L}[1]{>{\raggedright\arraybackslash}p{#1}}
  \caption{Summary of previous studies on Alzheimer's disease classification
    from neuroimaging. CN: cognitively normal; MCI: mild cognitive impairment;
    AD: Alzheimer's disease. Metrics are those reported by each study.}
  \label{tab:research_matrix}
  \begin{tabular}{L{2.2cm} L{2.4cm} L{1.5cm} L{2.0cm} L{3.9cm} L{3.7cm}}
    \toprule
    \textbf{Author (year)} & \textbf{Dataset} & \textbf{Modality} & \textbf{Task} &
    \textbf{Architecture} & \textbf{Validation metrics} \\
    \midrule
    % doi: 10.1016/j.neucom.2018.11.111 | note: Unsupervised; dataset not named in review
    Bi et al. (2020) & 1,075 subjects & MRI & AD vs CN & PCANet + K-means (unsupervised) & Acc 92.6\%; Prec 0.84; Rec 0.86 \\
    % doi: 10.1016/j.asoc.2019.105857 | note: Dataset not named in review
    Chitradevi \& Prabha (2020) & 200 subjects & MRI & AD vs CN & Grey Wolf Optimization (hippocampus sub-regions) + AlexNet & Acc 95\%; Prec 0.94; Rec 0.95 \\
    % doi: 10.1016/j.neucom.2020.01.053 | note: Precision/recall exceed accuracy: inconsistent in review
    Li et al. (2020) & ADNI (389 subjects) & 4D fMRI & AD vs CN & 3D CNN & Acc 88.12\%; Prec 0.92; Rec 0.93 \\
    % doi: 10.1016/j.neuroimage.2019.116459 | note: Segmentation + classification pipeline
    Liu et al. (2020) & ADNI (439 subjects) & MRI & CN vs MCI vs AD & Multi-model deep CNN + 3D DenseNet (hippocampus segmentation then classification) & Acc 76.2\%; Prec 0.78; Rec 0.77 \\
    % doi: 10.1016/j.patcog.2021.107944
    Chen \& Xia (2021) & ADNI (1,248 subjects) & MRI & CN vs MCI vs AD & Iterative sparse learning + deep learning & Acc 95.32\%; Prec 0.96; Rec 0.97 \\
    % doi: 10.1016/j.compbiomed.2021.104537
    Ebrahimi et al. (2021) & ADNI (450 subjects) & MRI & CN vs MCI vs AD & Temporal convolutional network (TCN) & Acc 91.78\%; Prec 0.87; Rec 0.90 \\
    % doi: 10.1016/j.bspc.2020.102397
    Hedayati et al. (2021) & ADNI (800 subjects) & MRI & CN vs MCI vs AD & Ensemble of convolutional autoencoders (feature extraction) & Acc 92.5\%; Prec 0.85; Rec 0.88 \\
    % doi: 10.1016/j.irbm.2020.06.006 | note: Modality as reported by review; verify
    Janghel \& Rathore (2021) & 2,729 images & fMRI, PET & AD vs CN & VGG16 + SVM & Acc 76\%; Prec 0.79; Rec 0.80 \\
    % doi: 10.1016/j.compbiomed.2021.104828 | note: Dataset not named in review
    Kaplan et al. (2021) & 263 subjects & MRI & CN vs MCI vs AD & LPQNet & Acc 99\%; Prec 0.93; Rec 0.94 \\
    % doi: 10.1016/j.cmpb.2021.106032 | note: Review and paper title disagree on architecture; review's DOI was wrong
    Liu et al. (2021) & OASIS (416 subjects) & sMRI & AD vs CN & Depthwise separable CNN & Acc 93\%; Prec 0.86; Rec 0.88 \\
    % doi: 10.1016/j.bspc.2021.103455
    Deepa \& Chokkalingam (2022) & ADNI (819 subjects) & MRI & CN vs MCI vs AD & VGG16 + Arithmetic Optimization Algorithm & Acc 97\%; Prec 0.99; Rec 0.98 \\
    % doi: 10.1007/s12559-021-09946-2
    Helaly et al. (2022) & ADNI (300 subjects) & MRI & CN vs MCI vs AD & CNN & Acc 98.47\%; Prec 0.95; Rec 0.96 \\
    % doi: 10.1016/j.bspc.2022.103565 | note: Precision/recall exceed accuracy: inconsistent in review
    Kong et al. (2022) & ADNI (740 subjects) & MRI + PET & AD vs CN & 3D CNN (multimodal) & Acc 86.52\%; Prec 0.94; Rec 0.95 \\
    % doi: 10.1016/j.matpr.2021.04.415 | note: Image-level count; check for subject-level split
    Kumar et al. (2022) & OASIS (5,000 images) & MRI & AD vs CN & AlexNet & Acc 97\%; Prec 0.93; Rec 0.92 \\
    % doi: 10.1016/j.asoc.2021.108099
    Sharma et al. (2022) & ADNI (400 subjects) & sMRI & CN vs MCI vs AD & FDN-ADNet: ResNet + fuzzy LS-TWSVM & Acc 97.29\%; Prec 0.98; Rec 0.98 \\
    % doi: 10.3389/frai.2025.1558725 | note: Public Kaggle slices; task differs between review and title
    Elmotelb et al. (2025) & Kaggle MRI (6,400 images) & MRI & AD vs CN; multi-class & ResNet152V2 + hyperparameter optimization & Acc 95\%; Prec 0.98; Rec 0.98 \\
    \bottomrule
  \end{tabular}
\end{table*}

% (requires the `array` and `booktabs` packages, both loaded in preamble.tex).
%
% Rows were compiled from the performance tables of Nibila & Sivagami (2026),
% so metrics are AS REPORTED BY THAT REVIEW and must be checked against each
% primary paper before the numbers are quoted. `%` comments above each row
% hold the DOI and verification notes; they are not typeset.
\begin{table*}[t]
  \centering
  \footnotesize
  \setlength{\tabcolsep}{3pt}
  \renewcommand{\arraystretch}{1.15}
  \newcolumntype{L}[1]{>{\raggedright\arraybackslash}p{#1}}
  \caption{Summary of previous studies on Alzheimer's disease classification
    from neuroimaging. CN: cognitively normal; MCI: mild cognitive impairment;
    AD: Alzheimer's disease. Metrics are those reported by each study.}
  \label{tab:research_matrix}
  \begin{tabular}{L{2.2cm} L{2.4cm} L{1.5cm} L{2.0cm} L{3.9cm} L{3.7cm}}
    \toprule
    \textbf{Author (year)} & \textbf{Dataset} & \textbf{Modality} & \textbf{Task} &
    \textbf{Architecture} & \textbf{Validation metrics} \\
    \midrule
    % doi: 10.1016/j.neucom.2018.11.111 | note: Unsupervised; dataset not named in review
    Bi et al. (2020) & 1,075 subjects & MRI & AD vs CN & PCANet + K-means (unsupervised) & Acc 92.6\%; Prec 0.84; Rec 0.86 \\
    % doi: 10.1016/j.asoc.2019.105857 | note: Dataset not named in review
    Chitradevi \& Prabha (2020) & 200 subjects & MRI & AD vs CN & Grey Wolf Optimization (hippocampus sub-regions) + AlexNet & Acc 95\%; Prec 0.94; Rec 0.95 \\
    % doi: 10.1016/j.neucom.2020.01.053 | note: Precision/recall exceed accuracy: inconsistent in review
    Li et al. (2020) & ADNI (389 subjects) & 4D fMRI & AD vs CN & 3D CNN & Acc 88.12\%; Prec 0.92; Rec 0.93 \\
    % doi: 10.1016/j.neuroimage.2019.116459 | note: Segmentation + classification pipeline
    Liu et al. (2020) & ADNI (439 subjects) & MRI & CN vs MCI vs AD & Multi-model deep CNN + 3D DenseNet (hippocampus segmentation then classification) & Acc 76.2\%; Prec 0.78; Rec 0.77 \\
    % doi: 10.1016/j.patcog.2021.107944
    Chen \& Xia (2021) & ADNI (1,248 subjects) & MRI & CN vs MCI vs AD & Iterative sparse learning + deep learning & Acc 95.32\%; Prec 0.96; Rec 0.97 \\
    % doi: 10.1016/j.compbiomed.2021.104537
    Ebrahimi et al. (2021) & ADNI (450 subjects) & MRI & CN vs MCI vs AD & Temporal convolutional network (TCN) & Acc 91.78\%; Prec 0.87; Rec 0.90 \\
    % doi: 10.1016/j.bspc.2020.102397
    Hedayati et al. (2021) & ADNI (800 subjects) & MRI & CN vs MCI vs AD & Ensemble of convolutional autoencoders (feature extraction) & Acc 92.5\%; Prec 0.85; Rec 0.88 \\
    % doi: 10.1016/j.irbm.2020.06.006 | note: Modality as reported by review; verify
    Janghel \& Rathore (2021) & 2,729 images & fMRI, PET & AD vs CN & VGG16 + SVM & Acc 76\%; Prec 0.79; Rec 0.80 \\
    % doi: 10.1016/j.compbiomed.2021.104828 | note: Dataset not named in review
    Kaplan et al. (2021) & 263 subjects & MRI & CN vs MCI vs AD & LPQNet & Acc 99\%; Prec 0.93; Rec 0.94 \\
    % doi: 10.1016/j.cmpb.2021.106032 | note: Review and paper title disagree on architecture; review's DOI was wrong
    Liu et al. (2021) & OASIS (416 subjects) & sMRI & AD vs CN & Depthwise separable CNN & Acc 93\%; Prec 0.86; Rec 0.88 \\
    % doi: 10.1016/j.bspc.2021.103455
    Deepa \& Chokkalingam (2022) & ADNI (819 subjects) & MRI & CN vs MCI vs AD & VGG16 + Arithmetic Optimization Algorithm & Acc 97\%; Prec 0.99; Rec 0.98 \\
    % doi: 10.1007/s12559-021-09946-2
    Helaly et al. (2022) & ADNI (300 subjects) & MRI & CN vs MCI vs AD & CNN & Acc 98.47\%; Prec 0.95; Rec 0.96 \\
    % doi: 10.1016/j.bspc.2022.103565 | note: Precision/recall exceed accuracy: inconsistent in review
    Kong et al. (2022) & ADNI (740 subjects) & MRI + PET & AD vs CN & 3D CNN (multimodal) & Acc 86.52\%; Prec 0.94; Rec 0.95 \\
    % doi: 10.1016/j.matpr.2021.04.415 | note: Image-level count; check for subject-level split
    Kumar et al. (2022) & OASIS (5,000 images) & MRI & AD vs CN & AlexNet & Acc 97\%; Prec 0.93; Rec 0.92 \\
    % doi: 10.1016/j.asoc.2021.108099
    Sharma et al. (2022) & ADNI (400 subjects) & sMRI & CN vs MCI vs AD & FDN-ADNet: ResNet + fuzzy LS-TWSVM & Acc 97.29\%; Prec 0.98; Rec 0.98 \\
    % doi: 10.3389/frai.2025.1558725 | note: Public Kaggle slices; task differs between review and title
    Elmotelb et al. (2025) & Kaggle MRI (6,400 images) & MRI & AD vs CN; multi-class & ResNet152V2 + hyperparameter optimization & Acc 95\%; Prec 0.98; Rec 0.98 \\
    \bottomrule
  \end{tabular}
\end{table*}


% ---------------------------------------------------------------------
% CITATION KEYS: the papers below are not yet in references.bib. Keys are
% PREDICTED from the Better BibTeX pattern (auth.lower + year +
% shorttitle(1,0)); after importing the DOIs in Zotero, replace any key
% that differs. DOI / key map:
%   liu2020multi        10.1016/j.neuroimage.2019.116459
%   liu2021alzheimers   10.1016/j.cmpb.2021.106032
%   kumar2022alexnet    10.1016/j.matpr.2021.04.415
%   hedayati2021deep    10.1016/j.bspc.2020.102397
%   kong2022multi       10.1016/j.bspc.2022.103565
%   chen2021iterative   10.1016/j.patcog.2021.107944
%   chitradevi2020analysis 10.1016/j.asoc.2019.105857
%   bi2020computer      10.1016/j.neucom.2018.11.111
%   ebrahimi2021deep    10.1016/j.compbiomed.2021.104537
%   kaplan2021feed      10.1016/j.compbiomed.2021.104828
%   sharma2022fdn       10.1016/j.asoc.2021.108099
%   elmotelb2025novel   10.3389/frai.2025.1558725
%   zhao2024vision      10.3389/fneur.2024.1490829   (read in full text)
%   mahmudjoy2025vitad  10.1016/j.brainres.2024.149302 (abstract only)
% ---------------------------------------------------------------------

\subsection{Convolutional architectures and transformers}
\label{sec:rw_architectures}

Convolutional neural networks (CNNs) dominate the neuroimaging studies in
\cref{tab:research_matrix}, but the architectures have moved well beyond plain
2D classifiers. \textcite{liu2020multi} chain a multi-model CNN for hippocampus
segmentation with a 3D DenseNet classifier, while \textcite{kong2022multi} apply
3D convolutions to fused MRI and PET volumes, and \textcite{ebrahimi2021deep}
treat the slices of a scan as a sequence processed by a temporal convolutional
network. Lighter and deeper backbones are also explored, from depthwise
separable convolutions \parencite{liu2021alzheimers} to a ResNet152V2 tuned by
hyperparameter optimization \parencite{elmotelb2025novel}.

Transformers appear only recently in this literature. \textcite{zhao2024vision}
argue that a pure 3D vision transformer is too costly for volumetric MRI and
instead redesign a 3D ResNet-50 with transformer-inspired choices: a
non-overlapping patchify stem, a Swin-like 1:1:3:1 block ratio, depthwise
separable convolutions, GELU activations and layer normalization. Using 818
ADNI subjects and a subject-level split, they report 92.14\% accuracy on the
three-class CN/MCI/AD task, although on the binary CN vs.\ AD task their
94.07\% remains below a previously published 98.20\%. At the other extreme, ViTAD, a
modified vision transformer, reports 99.98\%
accuracy on five classes using 1,296 ADNI images \parencite{mahmudjoy2025vitad}.
Results this close to perfect on a small, single-dataset image collection call
for scrutiny of how images were split across training and test sets, a point we
return to in \cref{sec:rw_caveats}.

\subsection{Support vector machines and their variants}
\label{sec:rw_svm}

Classical machine learning remains competitive when it is fed good features.
The review by \textcite{nibila2026artificial} collects SVM-based studies whose
reported accuracy spans from 69.5\% (ADNI, 1,715 subjects) to 95\% (150
subjects), and the lowest value belongs to the largest cohort, which suggests that
small samples may inflate reported performance. More recent work replaces
the standard SVM with margin-based variants and couples them to learned features:
\textcite{sharma2022fdn} feed ResNet features into a fuzzy least-squares twin SVM
(FDN-ADNet) and report 97.29\% accuracy on 400 ADNI subjects for the
three-class task.

\subsection{Feature extraction}
\label{sec:rw_features}

A second line of work concentrates on \emph{what} is given to the classifier.
Some studies restrict the input to regions known to be affected early in AD:
\textcite{chitradevi2020analysis} optimize the selection of hippocampal
sub-regions before an AlexNet classifier, and \textcite{liu2020multi}
segment the hippocampus before classifying it. Others learn the representation
without labels, as in the ensemble of convolutional autoencoders of
\textcite{hedayati2021deep} or the unsupervised PCANet plus $k$-means pipeline of
\textcite{bi2020computer}. A third group combines sparse learning with deep
networks \parencite{chen2021iterative} or rely on a feed-forward network built around local phase quantization
features, as in the LPQNet of \textcite{kaplan2021feed}. Unlike end-to-end CNNs,
these pipelines make the feature-extraction step an explicit design decision.

\subsection{Reporting practices and caveats}
\label{sec:rw_caveats}

Direct comparison across these studies is difficult. Cohort sizes range from a
few hundred subjects to several thousand images, and several studies count images
rather than subjects: for example, \textcite{kumar2022alexnet} report 97\%
accuracy for AlexNet on 5,000 OASIS images, and \textcite{elmotelb2025novel} use
6,400 public slices. When slices from one subject can fall into both training and
test sets, accuracy is overestimated, so a subject-level split is essential;
\textcite{zhao2024vision} state explicitly that they use one. Tasks also differ (CN vs.\ AD,
CN/MCI/AD, MCI vs.\ AD), as do the metrics reported, and
\textcite{nibila2026artificial} stress that accuracy alone is insufficient and
should be accompanied by precision, recall and related measures.
